\documentclass[11pt]{article}
\usepackage[margin=1in]{geometry}
\usepackage{amsmath,amssymb,amsthm}
\usepackage[hidelinks]{hyperref}
\newtheorem{theorem}{Theorem}
\newtheorem{lemma}[theorem]{Lemma}
\newtheorem{proposition}[theorem]{Proposition}
\newtheorem{corollary}[theorem]{Corollary}
\theoremstyle{remark}
\newtheorem{remark}[theorem]{Remark}
\newtheorem{example}[theorem]{Example}
\DeclareMathOperator{\Nil}{Nil}
\DeclareMathOperator{\Per}{Per}
\DeclareMathOperator{\rad}{rad}
\title{Additive periods of power maps lie in the prime radical}
\author{}
\date{}
\begin{document}
\maketitle
\begin{abstract}
We prove that every additive period of every power map on an associative ring
belongs to its prime radical. No identity, characteristic, finiteness, or
polynomial-identity assumption is required. The proof combines a substitution
using the highest nonzero power of a nilpotent element with the classical
bounded-index one-sided nil-ideal lemma. It follows that every nilperiod ring
is 2-primal and that its nilpotents form a locally nilpotent ideal. This
strengthens the NI conclusion conjectured by Burnette. We also characterize
finite unital nilperiod rings as finite products of local rings, and determine
the periodic power maps on Corbas rings, correcting the count for nonidentity
twists.
\end{abstract}

\section{Introduction}
All rings are associative and may lack an identity unless otherwise stated.
For a ring $R$, let $\Nil(R)$ be its set of nilpotent elements, and let $P(R)$
be its prime radical, also called the lower nilradical. Thus $P(R)$ is the
intersection of the prime ideals, or equivalently the smallest ideal for which
$R/P(R)$ is semiprime. The prime radical is locally nilpotent: every finite
subset generates a nilpotent subring. In particular, $P(R)$ is a nil ideal;
see, for example, \cite[Section~1]{Nielsen}.

For an integer $n\geq1$, set $\alpha_n(x)=x^n$ and define its additive period
group by
\[
 \Per(\alpha_n)=\{a\in R:(x+a)^n=x^n\text{ for every }x\in R\}.
\]
Evaluation at zero gives $\Per(\alpha_n)\subseteq\Nil(R)$. Our main result
locates these periods more precisely.

\begin{theorem}\label{thm:main}
For every associative ring $R$ and every positive integer $n$,
\[
 \Per(\alpha_n)\subseteq P(R).
\]
Consequently, every power map on a semiprime ring has trivial additive period
group.
\end{theorem}

Following Burnette \cite{Burnette}, a ring is \emph{nilperiod} if every
nilpotent element belongs to $\Per(\alpha_n)$ for some $n$, which may depend
on the element. A ring is \emph{NI} if its nilpotents form a two-sided ideal,
and it is \emph{2-primal} if $\Nil(R)=P(R)$.

\begin{corollary}\label{cor:nilperiod}
Every nilperiod ring is 2-primal. Its nilpotent elements form a locally
nilpotent two-sided ideal, which is both its prime radical and its largest nil
ideal. In particular, every nilperiod ring is an NI-ring.
\end{corollary}
\begin{proof}
Theorem~\ref{thm:main} gives $\Nil(R)\subseteq P(R)$, while the reverse
containment holds in every ring. The other conclusions follow from the
standard properties of the prime radical.
\end{proof}

Burnette's Conjecture~2.6 asks whether every nilperiod ring is NI, and his
Theorem~2.9 establishes this for PI algebras. Corollary~\ref{cor:nilperiod}
proves the conjecture and strengthens its conclusion. Theorem~\ref{thm:main}
also treats individual periods in rings that need not be nilperiod. In
particular, it extends the absence of nonzero periods in matrix rings over
division rings used in \cite[Lemma~2.8]{Burnette} to all semiprime rings.

The main proof uses the classical fact that a semiprime ring has no nonzero
one-sided nil ideal of bounded index. The new reduction is short: given a
period $a$ of nilpotency index $j$, evaluate its identity at
$x=a^{j-1}r$ and multiply by $a^{j-2}r$. This produces the bounded-index
right ideal to which the classical result applies.

The remaining results concern finite rings and power-map counts. As usual,
a unital ring is local if its nonunits form a proper ideal.
\begin{theorem}\label{thm:finite}
A nonzero finite unital ring is nilperiod if and only if it is a finite direct
product of finite local rings.
\end{theorem}
The zero ring may be included as the empty product. The structural
decomposition behind this result is standard; we explicitly credit it in
Section~3. Section~4 gives an exact computation for the twisted rings in
\cite[Example~2.2]{Burnette}.

\section{Periods in semiprime rings}
We use the following classical bounded-index lemma. Its left-ideal form is
stated in \cite[p.~93, item~(V)]{BM}; the right-ideal form follows by passing
to the opposite ring. The bounded nilradical discussion in
\cite[Section~1]{Nielsen} also records the result in the nonunital setting.

\begin{lemma}[Bounded-index one-sided nil ideals]\label{lem:levitzki}
Let $S$ be an associative ring, not necessarily unital. If $S$ has a nonzero
right ideal $I$ and a fixed positive integer $N$ such that $y^N=0$ for every
$y\in I$, then $S$ has a nonzero nilpotent two-sided ideal. In particular,
a semiprime ring has no nonzero nil right ideal of bounded index.
\end{lemma}

This lemma requires neither a chain condition on ideals nor a restriction on
the characteristic. The uniform bound concerns the elements of $I$, not all
nilpotents of $S$.

\begin{proof}[Proof of Theorem~\ref{thm:main}]
First let $S$ be semiprime and let $a\in\Per(\alpha_n)$. Since $a^n=0$,
a nonzero $a$ has a minimal nilpotency index $j$ satisfying $2\leq j\leq n$.
Suppose that this is the case, and put $b=a^{j-1}\ne0$.

Let $S^1=\mathbb Z\oplus S$ be the standard unitization, with multiplication
\[
 (m,u)(k,v)=(mk,mv+ku+uv).
\]
If $S$ already has an identity, we may take $S^1=S$. For any $r\in S^1$,
the element $x=br$ lies in $S$, so the period identity applies to $x$. Also
$ax=a^jr=0$. In the word expansion of $(x+a)^n$, a word vanishes whenever
an $a$ occurs before an $x$, since such a word contains adjacent factors $ax$.
The remaining words have the form $x^{n-k}a^k$, and terms with $k\geq j$
vanish. Hence
\begin{equation}\label{eq:expansion}
 0=(x+a)^n-x^n=\sum_{k=1}^{j-1}x^{n-k}a^k.
\end{equation}
Right-multiply \eqref{eq:expansion} by $a^{j-2}r$, interpreting $a^0=1$
in $S^1$ if $j=2$. The term with $k=1$ becomes
\[
 x^{n-1}a^{j-1}r=x^n,
\]
while all terms with $k\geq2$ vanish because $a^{k+j-2}=0$. Therefore
\[
 (br)^n=0\qquad(r\in S^1).
\]
The nonzero principal right ideal $bS^1=\mathbb Zb+bS$ of $S$ is thus nil
of bounded index $n$. This contradicts Lemma~\ref{lem:levitzki}, so $a=0$.
Notice that we have only evaluated the period identity at elements of $S$;
no period property of $S^1$ is assumed.

Now take an arbitrary ring $R$ and $a\in\Per(\alpha_n)$. Its period identity
descends to the quotient $R/P(R)$. Since this quotient is semiprime, the
first part of the proof shows that the image of $a$ is zero. Thus $a\in P(R)$.
\end{proof}

\begin{remark}
Local nilpotence in Corollary~\ref{cor:nilperiod} does not supply a uniform
nilpotency exponent for the entire ideal $\Nil(R)$. Nor does the proof
assert that each nilpotent element generates a nilpotent two-sided ideal.
It does imply that every such generated ideal is nil, since it lies in
$P(R)$.
\end{remark}

\section{Idempotents and finite rings}
\begin{proposition}\label{prop:central}
Every idempotent in a nilperiod ring is central.
\end{proposition}
\begin{proof}
Let $e^2=e$ and $r\in R$. The element $a=er-ere$ satisfies
$a^2=0$, $ea=a$, and $ae=0$. Thus $e+a$ is idempotent.
A period identity for $a$, evaluated at $e$, gives
\[
 e+a=(e+a)^n=e^n=e.
\]
Hence $er=ere$. Similarly, $b=re-ere$ is square-zero and $e+b$ is
idempotent, so the same argument gives $re=ere$. Therefore $er=re$.
\end{proof}

\begin{lemma}\label{lem:idempotent-power}
Every element of a finite ring has an idempotent positive power.
\end{lemma}
\begin{proof}
For $x$ in a finite ring there exist positive integers $a,b$ such that
$x^{a+b}=x^a$. Choose $n\geq a$ divisible by $b$. Eventual periodicity
of the powers then gives $x^{2n}=x^n$.
\end{proof}

The forward structural implication is an instance of the standard fact that
an Abelian semiperfect ring is a finite direct product of local rings
\cite[Proposition~2.6]{HLP}. Finite rings are semiperfect. We include an
elementary proof in this finite setting.

\begin{proof}[Proof of Theorem~\ref{thm:finite}]
Suppose first that $R$ is finite, unital, and nilperiod. By
Proposition~\ref{prop:central}, its idempotents are central.
Successively split the identity into orthogonal nonzero central idempotents
until each component contains no idempotents except zero and its identity.
The process terminates since $R$ is finite. It gives
\[
 R\cong R_1\times\cdots\times R_t.
\]
Each direct factor is nilperiod: a nilpotent in it lifts to a nilpotent
with zero coordinates in the other factors, and the corresponding period
identity restricts to that factor. Its nilpotents therefore form an ideal
by Corollary~\ref{cor:nilperiod}.

For $x\in R_i$, Lemma~\ref{lem:idempotent-power} gives $x^n=0$ or
$x^n=1_{R_i}$. Thus $x$ is either nilpotent or a unit. The nonunits of
$R_i$ are exactly its nilpotents, which form a proper ideal, so $R_i$ is
local.

Conversely, let $R$ be finite and local, with nonunit ideal $M$.
Its only idempotents are zero and one: if $e^2=e$, one of $e,1-e$ is a
unit, and $e(1-e)=0$ gives the conclusion.
For $x\in M$, an idempotent power remains in $M$, so it must be zero.
Thus $\Nil(R)=M$. Choose $n$ at least the maximum of the nilpotency
indices of the elements of $M$, and divisible by the exponent of the
finite group $R^\times$. Then
\[
 x^n=\begin{cases}0,&x\in M,\\1,&x\notin M.\end{cases}
\]
Translation by any element of $M$ preserves membership in $M$.
Consequently every element of $\Nil(R)$ is a period of $\alpha_n$.
For a finite direct product of finite local rings, choose one $n$
satisfying these requirements in all factors.
\end{proof}

\begin{example}
An NI-ring need not be nilperiod, even when finite. The upper triangular
$2\times2$ matrices over $\mathbb F_2$ have nilpotents
$\{0,E_{12}\}$, a square-zero ideal. But for $e=\operatorname{diag}(1,0)$
and $a=E_{12}$, the distinct matrices $e$ and $e+a$ are both idempotent.
Hence $(e+a)^n\ne e^n$ for every $n\geq1$.
\end{example}

\section{Power maps on Corbas rings}
Let $q=p^f$, where $p$ is prime and $f\geq1$, and choose $0\leq s<f$.
Let $\phi(a)=a^{p^s}$ be the corresponding automorphism of $\mathbb F_q$,
and put $g=\gcd(f,s)$. We allow $s=0$, so the identity automorphism has
$g=f$. Consider the Corbas ring
\begin{equation}\label{eq:corbas}
 R=\mathbb F_q\oplus\mathbb F_q,\qquad
 (a,b)(c,d)=(ac,ad+b\phi(c)).
\end{equation}
This is an associative unital ring, with identity $(1,0)$ and square-zero
ideal $N=\{0\}\oplus\mathbb F_q$. If $a\ne0$, the element $(a,b)$ has inverse
\[
 \bigl(a^{-1},-a^{-1}b\phi(a^{-1})\bigr).
\]
Thus $R$ is local and $\Nil(R)=N$.

The following calculation corrects the assertion in
\cite[Example~2.2 and Table~1]{Burnette} that nonidentity twists have no
periodic power maps. It also determines exactly which exponents occur.

\begin{theorem}\label{thm:corbas}
For the ring \eqref{eq:corbas}, set
\[
 M=\frac{q-1}{p^g-1}.
\]
Then
\[
 \Per(\alpha_n)=
 \begin{cases}
 N,&pM\mid n,\\
 \{0\},&pM\nmid n.
 \end{cases}
\]
There are exactly $p^g-1$ distinct periodic power maps on $R$.
\end{theorem}
\begin{proof}
Direct induction gives
\[
 (a,b)^n=\bigl(a^n,bC_n(a)\bigr),\qquad
 C_n(a)=\sum_{j=0}^{n-1}a^j\phi(a)^{n-1-j}.
\]
Every additive period must lie in $N$, by evaluation at zero. For $y\ne0$,
the element $(0,y)$ is a period exactly when $C_n(a)=0$ for all
$a\in\mathbb F_q$. At $a=1$ this forces $p\mid n$, and hence $n\geq2$.
The condition at $a=0$ then holds. If $a\ne0$ and $a=\phi(a)$, then
$C_n(a)=na^{n-1}$, which also vanishes. For $a\ne\phi(a)$ the geometric
sum gives
\[
 C_n(a)=\frac{a^n-\phi(a)^n}{a-\phi(a)}.
\]
Its vanishing is equivalent to $(a/\phi(a))^n=1$.
The homomorphism $a\mapsto a/\phi(a)$ on the cyclic group
$\mathbb F_q^\times$ has kernel $\mathbb F_{p^g}^\times$ and therefore
image of order $M$. Consequently $C_n$ vanishes everywhere if and only
if $p\mid n$ and $M\mid n$. Since $p\nmid M$, this is $pM\mid n$.
The argument also determines the full period group as stated.

For the count, the unit group of $R$ has exponent $E=p(q-1)$.
The displayed power formula shows that every unit has $E$th power one.
Elements $(a,0)$ realize order $q-1$, while $(1,b)$ with $b\ne0$ have
order $p$, so the exponent is exactly $E$.
All nonunits have square zero. Therefore, for $n,m\geq2$, the maps
$\alpha_n$ and $\alpha_m$ agree exactly when $n\equiv m\pmod E$:
sufficiency follows from the unit and nonunit cases, and necessity from
the definition of the unit-group exponent.
There are thus $E$ distinct maps with exponent at least two, and among
their exponent classes exactly
\[
 \frac{E}{pM}=p^g-1
\]
are periodic. The identity map $\alpha_1$ is distinct from these maps on
nonzero elements of $N$ and has no nonzero additive period.
\end{proof}

\begin{example}
For $q=4$ and the nontrivial automorphism $\phi(a)=a^2$, one has
$p=2$, $g=1$, $M=3$, and $E=6$. The power map $\alpha_6$ equals
$(1,0)$ on all units and zero on $N$, so its period group is $N$.
It is the unique distinct periodic power map on this noncommutative
sixteen-element ring.
\end{example}

\begin{thebibliography}{9}
\bibitem{BM}
H.~E.~Bell and W.~S.~Martindale III,
\emph{Centralizing mappings of semiprime rings},
Canad. Math. Bull. \textbf{30} (1987), no.~1, 92--101.
\href{https://doi.org/10.4153/CMB-1987-014-x}{doi:10.4153/CMB-1987-014-x}.

\bibitem{Burnette}
C.~Burnette,
\emph{On power maps over periodic rings},
Graduate J. Math. \textbf{10} (2025), no.~1, 5--16.
\href{https://gradmath.org/wp-content/uploads/2025/07/GJM2025-Burnette.pdf}{Published article}.
Earlier preprint: \emph{On power maps over weakly periodic rings},
\href{https://arxiv.org/abs/2207.14283}{arXiv:2207.14283}.

\bibitem{HLP}
J.~Han, Y.~Lee, and S.~Park,
\emph{Structure of Abelian rings},
Front. Math. China \textbf{12} (2017), no.~1, 117--134.
\href{https://doi.org/10.1007/s11464-016-0586-z}{doi:10.1007/s11464-016-0586-z}.

\bibitem{Nielsen}
P.~P.~Nielsen,
\emph{Bootstrapping the bounded nilradical},
J. Pure Appl. Algebra \textbf{217} (2013), no.~9, 1711--1715.
\href{https://doi.org/10.1016/j.jpaa.2012.12.002}{doi:10.1016/j.jpaa.2012.12.002}.
\end{thebibliography}
\end{document}
